\documentclass[11pt]{article}
\usepackage[margin=25mm]{geometry}
\usepackage{fontspec}
\setmainfont{Latin Modern Roman}
\usepackage{amsmath,amssymb,mathtools}
\usepackage{unicode-math}
\setmathfont{Latin Modern Math}
\usepackage{xurl,hyperref}
\hypersetup{colorlinks=true,urlcolor=blue}
\setcounter{secnumdepth}{0}
\setlength{\parindent}{0pt}
\setlength{\parskip}{5pt}
\setlength{\emergencystretch}{3em}
\newcommand{\sref}[2]{\hyperlink{source:#1:#2}{[S#2]}}
\newcommand{\cataloguescope}[1]{\paragraph{Recorded target.}#1}
\newcommand{\R}{\mathbb R}
\newcommand{\C}{\mathbb C}
\newcommand{\N}{\mathbb N}
\newcommand{\Z}{\mathbb Z}
\newcommand{\Q}{\mathbb Q}
\newcommand{\F}{\mathbb F}
\newcommand{\eps}{\varepsilon}
\begin{document}
% UnsolvedMath ID 30005664; source catalogue code OWR-14297742-004
\setcounter{section}{30005663}
\section{Aubin–Talenti Optimizers for Critical Dirac Potentials}\label{30005664}
{\small UnsolvedMath ID 30005664 · Mathematical Physics}
\subsection{Definitions and mathematical statement}

\paragraph{Shared notation.}$\mathbb N=\{1,2,\ldots\}$; the standard number fields have their usual meanings. All additional conventions are stated locally.


For fixed mass $m>0$, the massive Dirac operator $D_m$ has essential spectral gap $(-m,m)$. Take $d>1$ and $p\geq d$, and let $V\in L^p(\mathbb R^d)$ be nonnegative. Write $\lambda_D(V)$ for the lowest eigenvalue of the perturbed operator $D_m-V$ in the gap, in the admissible regime of the source, and $\Lambda_D(\alpha,p)=\inf_{\|V\|_p=\alpha}\lambda_D(V)$. The critical value $\alpha_*(p,d)$ is the endpoint where the optimized gap eigenvalue reaches $-m$ as an eigenvalue or resonance. The proposed Aubin--Talenti profile has positive scale constants $A,B$, and may be translated.

\paragraph{Catalogue wording retained for traceability.}
In dimensions $d>1$ at the critical parameter $\alpha=\alpha_*(p,d)$, is the optimizing nonnegative potential for the Dirac spectral-gap problem of Aubin–Talenti form $V(x)=A/(1+B|x|^2)$?

\paragraph{Scope and source qualification.}The admissible exponent range and critical parameter are those of the cited spectral problem, not arbitrary choices of $p$ or potential norm.
\subsection{Short English statement}
Determine whether the critical optimizing Dirac potential in every dimension greater than one has the proposed Aubin--Talenti form.
\subsection{Research attempt}
\paragraph{Model, date and investigated scope.}
OpenAI GPT-6 Astra (configured model \texttt{gpt-6-astra}), ultra reasoning; generated 11 October 2026. Computed the proposed potential normalization and a conditional interior Euler relation, identifying where the threshold invalidates that route.

\paragraph{Sources and literature check.}
Read the full supplied definition, catalogue wording, source qualifications and listed references for 30005664. No fresh status or novelty search was performed. The argument below proves its auxiliary statements directly and does not import an unverified external theorem.

\paragraph{Normalization and a necessary interior stationarity condition.} For the proposed profile \(V(x)=A/(1+B|x|^2)\), \(A,B>0\), change variables \(y=\sqrt Bx\). Then
\[
 \|V\|_p^p=A^p B^{-d/2}C_{d,p},\qquad
 C_{d,p}=|S^{d-1}|\int_0^\infty\frac{r^{d-1}}{(1+r^2)^p}\,dr.
\]
The integral is finite exactly when \(2p>d\), which includes the stated \(p\geq d>1\). Fixing the norm therefore fixes \(AB^{-d/(2p)}\), not the scale \(B\) separately.
\paragraph{Variation calculation.} Suppose, only for this test, that an interior optimizer has a normalized simple isolated gap eigenvector \(\psi\), that its eigenvalue is differentiable under compactly supported potential variations, and that the constraint admits a Lagrange multiplier \(\eta\). Differentiating \((D_m-V)\psi=\lambda\psi\), pairing with \(\psi\), and cancelling the eigenvector derivative gives
\[
 \delta\lambda=-\int|\psi|^2\delta V,
 \qquad |\psi|^2=\eta pV^{p-1}
\]
on regions of strictly positive \(V\). Hence the proposed profile would require \(|\psi(x)|^2\) proportional to \((1+B|x|^2)^{-(p-1)}\). For a nonzero normalized eigenvector this necessary tail is integrable only if \(2(p-1)>d\).
\paragraph{Threshold audit.} This does not disprove the proposal where that integrability test fails: the critical state may be a resonance, and at \(-m\) the assumed isolated-eigenvalue differentiability need not survive. The first gap is a justified threshold variational equation with the correct state space and compactness. The source optimizer, its existence and its exact profile for all admissible \((d,p)\) remain undetermined; the displayed interior calculation cannot be silently extended to the spectral edge.

\paragraph{Verification and adversarial audit.}
Checked radial integrability, scale dependence, the sign in eigenvalue variation and why normalized eigenvector assumptions fail at a resonance. Model self-audit only; no independent human mathematical verification. Applied contribution prompt A before drafting and B afterwards. Checked radial integrability, scale dependence, the sign in eigenvalue variation and why normalized eigenvector assumptions fail at a resonance. First remaining gap: Justify compactness and a threshold variational equation for the eigenvalue-or-resonance problem at alpha star. The inherited source status is preserved; original-target completion remains zero.

Coordinator checked the profile norm, conditional isolated-eigenvalue variation and integrability threshold, with resonances excluded from the inference. This additional check was model-assisted, not independent human verification.

\paragraph{Tools, process and contributions.}
Python 3 and local file tools used to read definitions and serialize manuscripts. Calculations are symbolic reasoning; no numerical experiment or theorem prover. Contribution prompt A was followed by prompt B. GPT-6 Astra with ultra reasoning drafted the research and performed its stated self-audit; any additional coordinator contribution is identified in the verification and audit. A GPT-6-based Codex coordinator assembled the manuscript and checked repository integration; its exact variant and reasoning configuration were not exposed. No human mathematical contribution, independent human proof check or formal proof certification is supplied. Complete derivations and the consulted source references are included here; private reasoning transcripts are not published.

\paragraph{Final status and first remaining gap.}
\textbf{UNRESOLVED}. Justify compactness and a threshold variational equation for the eigenvalue-or-resonance problem at alpha star.

\section{Completion Estimate}
\noindent\textbf{COMPLETION: 0\%}
This is an elementary auxiliary calculation or an explicit test of a restricted strategy. No new portion of the full source target is established.

\subsection{Sources}
\begin{itemize}
\item[\hypertarget{source:30005664:1}{S1}] ULAM.AI and the UnsolvedMath Contributors, UnsolvedMath, record 30005664 (OWR-14297742-004). Catalogue curation under CC BY 4.0; definitions and English exposition added locally.
\url{https://huggingface.co/datasets/ulamai/UnsolvedMath}.
\item[\hypertarget{source:30005664:2}{S2}] Many-Body Quantum Systems (2023), Oberwolfach Reports, DOI 10.4171/owr/2023/39, pp. 2247--2302. Original report and mathematical context.
\url{https://doi.org/10.4171/owr/2023/39}.
\end{itemize}
\label{end:30005664}
\end{document}
